\section{Normalized-determinant graphs and support complexity}
\label{sec:tpc31-row-entropy}

Let $\cM$ be the finite set of physical row integers.  It is contained
in an interval of length at most $C_{\mathscr D}Q$, and each integer in
$\cM$ labels at most one row.  For $c\ge1$ and
$\Omega\subset\Z\setminus\{0\}$, define the bipartite adjacency matrix
\begin{equation}
 A_{c,\Omega}(\alpha,\gamma)
 =\one_{\,(m_\alpha-m_\gamma)/c\in\Omega}.
 \label{eq:tpc31-adjacency}
\end{equation}
The indicator already forces $c\mid m_\alpha-m_\gamma$.  The exclusion
of zero is consistent with the off-diagonal mask.

\begin{definition}[Physical determinant support count]
\label{def:tpc31-entropy}
Choose once and for all a constant $B_{\mathscr D}$ larger than the
row-span constant.  Put
\begin{equation}
 \cH_c(\Omega)
 :=\#\left(
 \Omega\cap
 \left[-\frac{B_{\mathscr D}Q}{c},
        \frac{B_{\mathscr D}Q}{c}\right]
 \right).
 \label{eq:tpc31-entropy}
\end{equation}
We call this cardinality the determinant support count, or support
complexity, of $\Omega$ at content $c$.  Its logarithm would be the
zero-order (Hartley) support entropy; no Shannon or metric entropy is
used in the argument.  The estimates below are stated directly in
terms of the count $\cH_c$.
\end{definition}

\begin{lemma}[Two-sided degree bound]
\label{lem:tpc31-degree}
Every row and every column of $A_{c,\Omega}$ contains at most
$\cH_c(\Omega)$ nonzero entries.
\end{lemma}

\begin{proof}
Fix a left row integer $m$.  For each allowed $\delta$, the equation
\[
 m-n=c\delta
\]
determines $n=m-c\delta$ uniquely.  Since row integers determine row
labels uniquely, there is at most one right label for this value.
Only values in the physical difference range can occur, so the row
degree is at most \eqref{eq:tpc31-entropy}.  Interchanging $m$ and $n$
gives the same column bound.
\end{proof}

\begin{lemma}[Bipartite Schur bound]
\label{lem:tpc31-schur}
Let $A=(A_{uv})$ be a finite zero--one matrix with maximum row degree
$K_+$ and maximum column degree $K_-$.  Then, for arbitrary complex
vectors $x$ and $y$,
\begin{equation}
 \left|\sum_{u,v}x_u y_v A_{uv}\right|^2
 \le K_+K_-\|x\|_2^2\|y\|_2^2.
 \label{eq:tpc31-schur}
\end{equation}
\end{lemma}

\begin{proof}
For each $u$, Cauchy--Schwarz gives
\[
 \left|\sum_vA_{uv}y_v\right|^2
 \le K_+\sum_vA_{uv}|y_v|^2.
\]
Summing in $u$ and using the column-degree bound yields
\[
 \|Ay\|_2^2\le K_+K_-\|y\|_2^2.
\]
Apply Cauchy--Schwarz to $\langle x,A y\rangle$.
\end{proof}

Combining the preceding two lemmas with the physical coefficient
bound gives the row estimate used in every later section.

\begin{proposition}[Weighted determinant-support bound]
\label{prop:tpc31-row-entropy}
For every $c\ge1$ and $\Omega\subset\Z\setminus\{0\}$,
\begin{equation}
 \boxed{
 \sum_{\alpha,\gamma}
 |\gamma_\alpha^{(1)}\gamma_\gamma^{(2)}|
 A_{c,\Omega}(\alpha,\gamma)
 \ll_{\mathscr D}
 Q\log(2L)\,\cH_c(\Omega).}
 \label{eq:tpc31-row-entropy-bound}
\end{equation}
\end{proposition}

\begin{proof}
Apply \cref{lem:tpc31-schur} to the nonnegative vectors
$|\gamma^{(1)}|$ and $|\gamma^{(2)}|$, use
\cref{lem:tpc31-degree}, and then use
\eqref{eq:tpc31-row-l2} for the two norms.
\end{proof}

\subsection{Three useful support-count calculations}

The first calculation is independent of the location of a window.

\begin{lemma}[Translated windows and residue cells]
\label{lem:tpc31-window-residue-entropy}
Let $I\subset\R$ be any interval of length at most $W\ge0$.
Then
\begin{align}
 \cH_c(I\cap\Z)&\le W+1,
 \label{eq:tpc31-window-entropy}\\
 \cH_c(\{\delta\in I\cap\Z:\delta\equiv a\!\pmod q\})
 &\le 1+\frac Wq
 \label{eq:tpc31-window-residue-entropy}
\end{align}
for every $q\ge1$ and $a\pmod q$.
\end{lemma}

\begin{proof}
An interval of length $W$ contains at most $W+1$ integers.  Integers
in one class modulo $q$ are spaced by $q$, so their number in the
same interval is at most $1+W/q$.  Intersecting with the physical
range can only reduce either count.
\end{proof}

For a residue class without an imposed window, the physical range
retains the content dependence.

\begin{lemma}[Full-range residue support]
\label{lem:tpc31-full-residue-entropy}
Uniformly for $q\ge1$ and $a\pmod q$,
\begin{equation}
 \boxed{
 \cH_c(\{\delta\in\Z:\delta\equiv a\!\pmod q\})
 \ll_{\mathscr D}1+\frac Q{cq}.}
 \label{eq:tpc31-full-residue-entropy}
\end{equation}
No coprimality condition between $c$ and $q$ is needed.
\end{lemma}

\begin{proof}
The physical normalized-determinant interval has length
$O_{\mathscr D}(Q/c)$.  One class modulo $q$ occupies it at most
$O_{\mathscr D}(1+Q/(cq))$ times.
\end{proof}

\begin{remark}[The endpoint one cannot be removed]
\label{rem:tpc31-plus-one}
If $q$ or $c$ is larger than the relevant span, a nonempty residue
cell may still contain one normalized determinant.  The $+1$ terms in
\eqref{eq:tpc31-window-residue-entropy} and
\eqref{eq:tpc31-full-residue-entropy} are therefore genuine floors of
the deterministic graph estimate.
\end{remark}

\begin{remark}[Macroscopic translation]
\label{rem:tpc31-macroscopic-window}
Nothing in \cref{lem:tpc31-window-residue-entropy} depends on the
distance from $I$ to zero.  One may take
$I\subset[\kappa Q/c,2\kappa Q/c]$ for a fixed $\kappa>0$ whenever
that interval meets the physical difference range.  The resulting
estimate is not another deletion of near-diagonal rows.
\end{remark}
